\section{Three Failure Modes}
\label{sec:phenomenon}

\paragraph{(i) Transition blindness.} Paths crossing a semantic boundary often
induce no model transition at all. Among single-turn paths discovered by
random edits, $108/128 = 84.4\%$ induce zero model transition; on directly
targeted fresh paths the miss rate is $28/64 = 43.8\%$ (95\% CI
$[0.32,0.56]$), with flipmine at $31.3\%$. Both numbers are true; their
denominators differ. The random-edit screen over-samples rare tangential,
grazing crossings---precisely the shape the model misses most---while the
targeted sweep drives decisively across the boundary. The miss rate is
therefore not a fixed percentage but depends on how a change approaches the
semantic boundary. We report both, never averaged, always with denominators.

\paragraph{(ii) Compositional blindness (headline).} When \emph{both} atomic
edits of a scene are judged correctly, their joint compositional change is
still missed in $151/176 = 85.8\%$ of fresh-holdout cases (discovery
$62/66 = 93.9\%$, same direction). Matched single-edit controls miss far less
($61.0\%$ clean), so the failure is specific to the composition, not to the
edits. A model can know each part and still not know the change.

\paragraph{(iii) Consequential blindness.} The misses propagate. Fresh-holdout
event detection shows matched misses; action selection under a fixed action
library is invalid in $75.0\%$ of hard fresh cases ($\lambda{=}2$); the
decomposition (C1/C0) shows a C1 miss of $52.4\%$ against a C0 false-alarm
rate of $26.2\%$. Decision-interface analysis (U02) shows flipmine's gain is
in feasible-set coverage (refusal $0.43 \to 0.12$) rather than non-refusal
precision, and that the gain survives across scoring rules
(see the evidence table, reports/FINAL_EVIDENCE_TABLE.md).
